% !TEX root = ./../main.tex

\section{Related Work}
\label{sec:related-work}

We now explore existing proposals related to our work in three directions. Approaches that use type 
systems to verify \ac{COP} programs, existing verification approaches that use Refinement Types
(\eg in Reactive Programming, a paradigm used to implement \ac{COP}~\cite{watanabe14,leger23}),
and verification techniques in domains closely related to \ac{COP} such as Self-Adaptive Systems
and Actor Systems~\cite{watanabe14}.


%%
\subsection{Use of Type Systems in \acl{COP}}
\label{sec:types-cop}

There exist several Type-based solutions to verify the properties of \ac{COP} programs.
~\citet{aotani15} propose an abstraction to guarantee that layers relying on an imperative activation 
mechanism do not call methods on other layers that are inactive. This is done by modeling layer 
activation with Typestates. A Typestate denotes the run-time state of a resource so that we can 
determine when an operation can be applied to the resource, and how it changes the resource's state.
In particular, there are two Typestates in the approach. The type $\uparrow L$, denotes that layer 
$L$ is active, and the type $\downarrow L$, denotes that layer $L$ is inactive. This way of tracking the 
state of layers allows us to define the \scode{activate} and \scode{deactivate} operations in a way that 
they can only be applied to a consistent state, \ie the \scode{activate} function can only be applied to  
layers with type $\downarrow L$, and the \scode{deactivate} function can only be applied to layers with 
type $\uparrow L$.
The problem solved with Typestates focuses only on the consistency of activation/deactivation 
operations, and not on the interaction between layers. Using Refinement Types we can have the same 
guarantees when we add the \scode{isActive} predicate as a precondition to the \scode{run} function.

JCop~\cite{inoue14} has Type-level constructions to verify layer interactions, by explicitly stating the 
dependency between two layers in their type definition. For instance, given two layers \scode{L1} and 
\scode{L2}, the programmer can specify a dependency between them in JCop by using the 
\scode{requires} keyword, ensuring that \scode{L1} is always activated before \scode{L2} 
(\scode{public layer L1 requires L2}).
This approach, however, only works well for imperative activation mechanisms as it was briefly 
discussed in \fref{sec:cop}. No alternatives have been found to statically verify layer interactions with 
declarative activation mechanism, and our approach with Refinement Types puts forward a viable 
solution to do so.


%%
\subsection{Verification of Reactive Programs Using Refinement Types}
\label{sec:verification-reactive-programming}

in the \ac{RP} paradigm programs \emph{react} to values that vary through time.
It is common to find applications that require software adaptability (such as IoT systems or web 
applications) implemented with \ac{RP}. For example, Bellevue relies on an abstraction that dictates how 
the program should react to changes in the state (the \scode{DrawingModel}). The \scode{update} 
function, in particular, is the main source of the relevant time-varying values (state and messages).

Software verification is used in \ac{RP} to have a precise description of the desired run-time behavior of 
a reactive program. The work by~\citet{chen24} is an example of using Refinement Types to verify \ac{RP} programs. Refinement Types are used 
to describe post-conditions for functions that process time-varying data (\ie streams), as well as to 
declare invariants over such data. A type system is developed extending the Refinement Types 
language with predicates drawn from Linear Temporal Logic. These predicates are used to reason 
about the behavior of data streams. In particular, the type system introduces the 
$\{ w : b \; | \; \square \; \phi \}$ refinement to specify that $\phi$ must hold for all values of the stream 
$w$ from the current instant of time onwards, and the $\{ w : b \; | \; \bigcirc \; \phi \}$ refinement to 
describe that $\phi$ must hold in the next time instant. The resulting system allows programmers to 
formally verify, simulate and deploy Cyber-Physical Systems.

The use of Refinement Types to verify Cyber-Physical Systems is closely related to our proposal. First, 
Cyber-Physical Systems are an application domain of \ac{COP}, often benefiting from dynamic 
adaptation. Second, as mentioned before, there are \ac{COP} languages based on \ac{RP}, so the use 
of Refinement Types in an \ac{RP} context can translate naturally to reactive-based \ac{COP} 
languages. 

Moreover over, we note that our approach could benefit from this Temporal Logic extension to prove 
more sophisticated properties. For example, to prove that the \scode{ColorFillAction} always receives 
the canvas bitmap if, at some previous point in time, the user selected the \scode{ColorFill} tool, is an 
specification that could be written using Temporal Logic.


%%
\subsection{Verification in Actor Systems for Software Adaptability}
\label{sec:verification-actors}

EnsembleS~\cite{harvey21} is an actor-based language based on Multiparty Session Types to 
define an verify programs that can \emph{adapt} their behavior in response to the operating 
environment. \emph{Adaptability} in this context refers to the ability of sensing the environment and 
respond by \emph{discovering} new components, \emph{replacing} behavior or \emph{communicating} 
information. EnsembleS defines actors as single-threaded entities that do not share state, and have a 
channel-based message communication between each other. The programmer is able to create static 
topologies, where all participants are present throughout the duration of a session, as well as dynamic 
topologies where participants can connect or disconnect at a given point in the session.

The main contribution of this work is the static verification of dynamic interaction and replacement of 
functionality, which cannot be guaranteed by using Multiparty Session Types alone. EnsembleS offers 
the \scode{link} and \scode{unlink} statements to enable dynamic connectivity,  and verifies that the 
replacement of functionality does not violate the communication protocol that is already established by 
the Session Type. 
It is possible to consider an alternative formulation to verify \ac{COP} programs with Session Types, 
where one can granularly specify the sequence of steps of a program with Session Types, and model 
(de)activation with a replacement mechanism similar to the one defined for EnsambleS. This would 
greatly differ from the Refinement Types alternative, as the programmer would need to explicitly 
specify all steps in the flow of the \ac{COP} program. However, such an approach could suffer from 
the state explosion problem of interactions between adaptations.


%%
\subsection{Verification in Self-Adaptive Systems}
\label{sec:verification-as}

\citet{weyns12} present a survey of the existing formal methods for Self-Adaptive Systems. The 
survey compiles various studies from conferences and journals, and discusses the main methods that 
have been used to formally reason about such systems, as well as the main program properties of 
interest to the researchers.
The main modeling language for property specification is algebra, but a very limited set of studies 
actually used automated reasoning to evidence the validity of such properties. The ones that do, often 
use Model Checking or Theorem Proving to a lesser extent. Liveness, safety, and reachability are 
some of the common properties of interest, but most studies focus on verifying behavior specific to the 
use case or area of self-adaptation.
One of the main remarks from the survey, is that there are not many studies that make use of software 
verification in this area, and that there is a need of light-weight and pluggable verification 
tools~\cite{weyns12}. As validated from our approach, Refinement Types could be a good candidate in 
that regard, as they provide a seamless and pluggable integration of property specification within the 
same language in which the application is being developed.
Furthermore, Refinement Types shares some similarities with model checking in the sense that 
Refinement Types use \ac{SMT} solvers to validate the verification constraints, with model checkers 
using SAT solvers to check the satisfiability of a propositional formula that represents the program 
traces of the finite state system under verification~\cite{armando06}.


\endinput

